% ------------------------------------------------------------
% Chapter 7: Dark Matter and Dark Energy in RSVP
% ------------------------------------------------------------

\chapter{Dark Matter and Dark Energy in RSVP}
\label{chap:dark-sector}

\section{Overview}

A major motivation for the \rsvpabbr{} cosmology is the possibility of
explaining apparent dark--matter and dark--energy phenomena as emergent
consequences of lamphrodynamic entropy and flux, without introducing
new fundamental particles or fields beyond $(\PhiField,\vField,\SField)$.
This chapter develops this idea quantitatively, focusing on galaxy rotation
curves, cluster mass profiles, weak lensing, and cosmic acceleration.

\section{Lamphrodynamic Contributions to Effective Density}

Recall from \cref{sec:theta-flrw} in \cref{chap:friedmann-like} that the
lamphrodynamic energy density $\rho_{\mathrm{lam}}$ arises from
entropy production and lamphrodynamic couplings. In homogeneous
cosmology, $\rho_{\mathrm{lam}}$ may drive cosmic acceleration with
effective equation of state $w_{\mathrm{lam}}\approx -1$; in
inhomogeneous settings, it can mimic cold dark matter by enhancing
gravitational attraction.

\begin{remark}[Unification possibility]
A single lamphrodynamic field $\SField$ may account for both
dark--energy–like acceleration and dark--matter–like clustering,
depending on the regime and scale.
\end{remark}

\section{Galaxy Rotation Curves}

Consider a spherically symmetric galaxy with baryonic mass $M(r)$.
In Newtonian gravity, the circular velocity satisfies
\[
  v^2(r) = \frac{GM(r)}{r}.
\]
In \rsvpabbr{}, lamphrodynamic contributions modify the effective potential
$\Phi_{\mathrm{eff}}$, yielding
\[
  v^2_{\mathrm{RSVP}}(r)
  = \frac{G}{r}\bigl[M(r) + M_{\mathrm{lam}}(r)\bigr],
\]
where $M_{\mathrm{lam}}(r)$ is the effective lamphrodynamic mass enclosed
within radius $r$, defined by
\[
  M_{\mathrm{lam}}(r)
  := \int_0^r 4\pi s^2\,\rho_{\mathrm{lam}}(s)\,ds.
\]
Flat rotation curves at large $r$ arise naturally if $\rho_{\mathrm{lam}}$
falls slower than $1/r^2$ or approaches a constant asymptotically.

\section{Cluster Mass Profiles}

Galaxy clusters require substantial dark matter in \LCDM{} for dynamical
stability. In \rsvpabbr{}, a lamphrodynamic energy density
$\rho_{\mathrm{lam}}$ adding to $\rho$ may account for the inferred total
mass without new particles. X--ray and gravitational lensing measurements
of cluster profiles place strong constraints on the radial dependence of
$\rho_{\mathrm{lam}}$; viable lamphrodynamic models must reproduce observed
mass distributions.

\section{Weak Lensing}

Weak lensing probes the projected mass distribution along the line of
sight. Lamphrodynamic energy density contributes to the lensing
convergence via
\[
  \kappa_{\mathrm{RSVP}}
  \sim \int_{\text{LOS}}
       (\rho + \rho_{\mathrm{lam}})\,ds.
\]
Matching observed shear maps therefore constrains the combination
$\rho_{\mathrm{lam}} + \rho$ rather than $\rho$ alone. Current weak--lensing
surveys may already limit $\rho_{\mathrm{lam}}$ in certain regimes; detailed
phenomenology remains to be carried out.

\section{Cosmic Acceleration Without $\Lambda$}

From the effective Friedmann equation \eqref{eq:friedmann1-rsvp}, cosmic
acceleration occurs when
\[
  \rho_{\mathrm{lam}} + \rho + \frac{\Lambda}{8\pi G}
  \text{ is sufficiently large}
\]
with effective negative pressure $p_{\mathrm{lam}}\approx -\rho_{\mathrm{lam}}$.
If $\Lambda=0$, lamphrodynamic entropy alone may drive acceleration, provided
$\dot{\SField}$ is large enough at late times.

\section{Interpretation}

Lamphrodynamic entropy and flux can mimic both dark matter and dark energy,
potentially unifying these dark components in a single framework. Whether
the \rsvpabbr{} predictions match observed rotation curves, cluster dynamics,
and cosmological acceleration quantitatively depends on the detailed time and
space dependence of $\SField$ and its couplings $\alpha_i,\beta_i$.

\section{Discussion}

While \rsvpabbr{} offers an appealing conceptual unification of the dark
sector, significant work remains to fully develop and test lamphrodynamic
phenomenology against observational data. In subsequent chapters (notably
\cref{chap:cmb,chap:lss,chap:supernovae}) we confront \rsvpabbr{} with
observations, beginning with the cosmic microwave background.

\clearpage

